% operator-stability-v1 — publication template (writeup-aligned)
% Slots filled by scripts/generate_operator_paper.py
\documentclass[12pt]{article}
\usepackage[letterpaper,top=2cm,bottom=2cm,left=3cm,right=3cm]{geometry}
\usepackage[T1]{fontenc}
\usepackage[utf8]{inputenc}
\usepackage{lmodern}
\usepackage{microtype}
\usepackage{amsmath,amssymb,amsthm,mathtools}
\usepackage{booktabs}
\usepackage{xcolor}
\usepackage{enumitem}
\usepackage{xurl}
\usepackage[colorlinks=true,linkcolor=blue,urlcolor=blue,citecolor=blue]{hyperref}

\Urlmuskip=0mu plus 1mu
\urlstyle{tt}
\emergencystretch=2em

\setlength{\parskip}{0.4em}
\setlength{\parindent}{1.2em}
\setlist[itemize]{leftmargin=1.4em,itemsep=0.25em,topsep=0.35em}
\setlist[enumerate]{leftmargin=1.4em,itemsep=0.25em,topsep=0.35em}

\theoremstyle{plain}
\newtheorem{theorem}{Theorem}
\newtheorem{lemma}[theorem]{Lemma}
\theoremstyle{definition}
\newtheorem{definition}[theorem]{Definition}
\newtheorem{example}[theorem]{Example}
\theoremstyle{remark}
\newtheorem{remark}[theorem]{Remark}

% Breakable monospace for Lean paths / identifiers (xurl \path)
\newcommand{\leanpath}[1]{%
  {\raggedright\urlstyle{tt}\path{#1}}}


\definecolor{fundbox}{RGB}{245,245,248}
\definecolor{fundrule}{RGB}{180,180,190}

\title{Interval-Membership Preservation under Bounded Score Perturbations}
\author{Research Architecture Runtime\\[0.25em]
{\normalsize\texttt{operator/interval-membership}}}
\date{2026-07-29}

\begin{document}
\maketitle

\begin{abstract}
\leanpath{interval-membership} is a scalar region indicator preserved outside an $\varepsilon$-buffer under $|x'-x|\le\varepsilon$. Lean \texttt{LEAN\_FULL}.
\end{abstract}

\noindent
\begin{center}
\fcolorbox{fundrule}{fundbox}{%
  \begin{minipage}{0.92\linewidth}
  \centering
  \textbf{This theorem is:} \texttt{Primitive}
  \end{minipage}}
\end{center}
\vspace{0.6em}

\section{Problem}
{\raggedright\sloppy
The \leanpath{interval-membership} operator maps a scalar score to a discrete region label/bit.
\par}

\section{Stability notion}
Additive $|x'-x|\le\varepsilon$; outputs agree when $x$ is at least $\varepsilon$ inside its region.

\section{Definitions}
Region boundaries are thresholds (sign at $0$, abs-threshold at $T\ge 0$, interval $[L,U]$). Assumptions follow the Lean module.

\section{Theorem}
\begin{theorem}\label{thm:inv}
Under the module's buffer hypotheses, $|x'-x|\le\varepsilon$ preserves the operator output; sharpness pushes across a boundary when the buffer fails.
\end{theorem}

\section{Intuition}
Same one-dimensional buffer logic as thresholding, applied to the operator's cut set.

\section{Examples}
\begin{example}
For sign: $x=2$, $\varepsilon=0.5$ stays positive. For abs-threshold $T=3$, $x=5$, $\varepsilon=0.2$ stays pass.
\end{example}

\section{Proof}
Write $|x'-x|\le\varepsilon$. On the pass side of each cut, $x$ is at least $\varepsilon$ beyond the cut, so $x'$ cannot cross it; on the fail side, $x'$ cannot reach the cut. Sharpness uses a $\pm\varepsilon$ push from a point within $\varepsilon$ of the boundary (Theorem~\ref{thm:inv}), matching the Lean cases.

\section{Formal statement}
{\raggedright\sloppy
\begin{description}[style=nextline,leftmargin=1.1em,font=\normalfont\bfseries]
\item[Lean module] \leanpath{Research.Operators.IntervalMembership.Preservation}
\item[Source file] \leanpath{lean/Research/Operators/IntervalMembership/Preservation.lean}
\item[Theorems]\leavevmode\par\vspace{0.15em}\begin{itemize}[leftmargin=1.2em]
  \item \leanpath{interval\_membership\_preservation}
  \item \leanpath{interval\_membership\_sharpness}
\end{itemize}
\item[Certificate] \leanpath{lean/certificates/interval-membership/interval-membership-preservation}
\item[Status] \texttt{LEAN\_FULL} on Mathlib $\mathbb{R}$ (\texttt{REAL\_MATHLIB})
\end{description}
\par}

\section{Proof dependencies}
{\raggedright\sloppy
\begin{itemize}\item Interval arithmetic on $\mathbb{R}$.\item Threshold-style buffers.\end{itemize}
\par}

\section{Consequences}
{\raggedright\sloppy
Scalar building block; related: thresholding, multi-threshold, absolute-value-threshold.
\par}

\end{document}
